\section{The Unified PLR Sensitivity Pipeline}
\label{sec:pipeline}

All three channels (Sections~\ref{sec:dme}--\ref{sec:migdal}) and the
ionization model (Section~\ref{sec:ionization}) feed the \textbf{same}
downstream chain, implemented once in C++ and dispatched by
\texttt{model.type} in the JSON config (\texttt{dm\_electron},
\texttt{dark\_photon}, \texttt{migdal}):

\begin{center}
\begin{tabular}{@{}c@{}}
Rate CSV [events/kg/yr/eV] \quad (continuous $dR/dE$, or boxcar for dark photon) \\
$\big\downarrow$ \\
\texttt{RateTable::MakeTH1D} — bin-center interpolation onto uniform $E$
  histogram \\
(needs fine \texttt{nbins} for narrow boxcar features — Section~\ref{sec:binning}) \\
$\big\downarrow$ \\
\texttt{ChargeIonization::FoldToNe} — $\times$ exposure ($M\cdot T$), fold
  against $P(n_e\mid E)$ \\
$\big\downarrow$ \\
$n_e$ spectrum $S(n_e)$ \\
$\big\downarrow$ \\
\texttt{EfficiencyMC} / \texttt{PatternRates} — diffusion, readout noise,
  pattern classification \\
(\texttt{observable\_bins}: \texttt{"ne"} or \texttt{"pattern"}) \\
$\big\downarrow$ \\
Profile likelihood ratio scan — background model (\texttt{dc\_flat\_migration}
  or \texttt{Bp\_theta\_Br}) \\
$\big\downarrow$ \\
Upper limit at 90\% CL: \ $\sigmae(m_\chidm)$ (DM-e), \ $\kappa(m_{A'})$ (dark
  photon), \ $\sigman(m_\chidm)$ (Migdal)
\end{tabular}
\end{center}

\subsection{What differs between channels, and what does not}

\begin{center}
\begin{tabular}{@{}p{3.2cm}p{3.5cm}p{3.5cm}p{3.5cm}@{}}
\toprule
Stage & DM-electron & Dark photon & Migdal \\
\midrule
Rate engine & QEDark / QCDark2 & DarkELF \texttt{absorption} & DarkELF \texttt{Migdal} \\
Input spectrum shape & Continuous $dR/dE(E)$ & Monochromatic boxcar at $m_{A'}$ & Continuous $dR/d\omega$, peaked near $\Egap$ \\
Scan axis ($x$) & $m_\chidm$ & $m_{A'}$ & $m_\chidm$ \\
Scan axis ($y$, linear in) & $\sigmae$ & $\kappa^2$ & $\sigman$ \\
\texttt{ChargeIonization::FoldToNe} & identical call & identical call & identical call \\
\texttt{EfficiencyMC} / \texttt{PatternRates} & identical & identical & identical \\
PLR machinery & identical & identical & identical \\
Background model & identical & identical & identical \\
\bottomrule
\end{tabular}
\end{center}

Because the $y$-axis enters each rate \textbf{linearly} ($\sigmae$ in
Section~\ref{sec:dme}, $\kappa^2$ in Section~\ref{sec:darkphoton},
$\sigman$ in Section~\ref{sec:migdal}), all three channels reuse the same
fast-rescaling trick when regenerating grids: compute once at a reference
coupling, then scale analytically for every other grid point — this is
exact by construction, not an approximation.

\subsection{Config field summary (\texttt{model.type})}

\begin{verbatim}
// DM-electron
"model": { "type": "dm_electron", "mediator": "heavy|light",
           "rates_dir": "data/qcdark2_rates/...", ... }

// Dark photon
"model": { "type": "dark_photon", "rates_dir": "data/darkphoton_rates/...", ... }

// Migdal
"model": { "type": "migdal", "mediator": "heavy|light",
           "target_nucleus": "Si28", "nuclear_A": 28, "nuclear_Z": 14,
           "rates_dir": "data/migdal_rates/...", ... }
\end{verbatim}

\begin{verbatim}
// Shared ionization block — used identically by all three
"response": {
  "charge_ionization": {
    "table_csv":   "data/p100K_gap0p34_eh1p7778.csv",
    "band_gap_eV": 0.34,
    "eh_pair_eV":  1.7778,
    "scenario":    "Klein"
  }
}
\end{verbatim}
